% ============================================================
% 03-izvedenice-i-alternative.tex — PRVA VERZIJA
% ============================================================

Temeljni SMOTE algoritam, unato\v{c} svojoj popularnosti, ima nekoliko
uo\v{c}enih nedostataka --- od prekomjerne generalizacije do generiranja
sinteti\v{c}kih primjera u podru\v{c}jima ve\'{c}inske klase
\cite{fernandez2018}. Ovi nedostaci motivirali su razvoj brojnih
izvedenica koje nastoje adresirati pojedine probleme. U ovom poglavlju
bit \'{c}e opisane klju\v{c}ne izvedenice SMOTE algoritma, uklju\v{c}uju\'{c}i
geometrijska pro\v{s}irenja, te alternativni pristupi rje\v{s}avanju
problema neuravnote\v{z}enosti klasa.

\subsection{Borderline-SMOTE}

Borderline-SMOTE \cite{han2005} fokusira se na grani\v{c}ne primjere
manjinske klase koji su klju\v{c}ni za odre\dj{}ivanje granice odluke.
Opisat \'{c}e se kategorizacija manjinskih primjera u tri skupine na
temelju broja ve\'{c}inskih susjeda me\dj{}u $m$ najbli\v{z}ih: sigurni
(Safe), grani\v{c}ni (Danger) i \v{s}um (Noise). Zatim se opisuju
dvije varijante: Borderline-SMOTE1 koja preuzorkuje samo Danger primjere
koriste\'{c}i isklju\v{c}ivo manjinske susjede, i Borderline-SMOTE2 koja
dopu\v{s}ta i ve\'{c}inske susjede uz faktor $\lambda \in (0, 0.5)$.
Dat \'{c}e se pseudo-k\^{o}d za obje varijante. Prednost je fokusiranje
na grani\v{c}ne primjere, a nedostatak osjetljivost na odabir parametra
$m$.

\subsection{ADASYN}

ADASYN (Adaptive Synthetic Sampling) \cite{he2008adasyn} uvodi adaptivni
pristup: vi\v{s}e sinteti\v{c}kih primjera generira se u podru\v{c}jima
gdje je ve\'{c}inska klasa gu\v{s}\'{c}a u susjedstvu manjinskih primjera.
Opisat \'{c}e se ra\v{c}unanje omjera $\Gamma_i = \Delta_i / k$ gdje je
$\Delta_i$ broj ve\'{c}inskih susjeda me\dj{}u $k$ najbli\v{z}ih,
normalizacija u distribuciju $\hat{\Gamma}_i$ te raspodjela ukupnog broja
sinteti\v{c}kih primjera $G$ prema toj distribuciji. Dat \'{c}e se
pseudo-k\^{o}d. Automatsko fokusiranje na te\v{s}ke primjere je glavna
prednost, dok je mana osjetljivost na outlier-e koji dobivaju najvi\v{s}e
sinteti\v{c}kih primjera.

\subsection{Safe-Level-SMOTE}

Safe-Level-SMOTE \cite{bunkhumpornpat2009} uvodi koncept sigurnosne razine
(safe-level) za svaki manjinski primjer --- broj manjinskih susjeda me\dj{}u
$k$ najbli\v{z}ih. Opisat \'{c}e se strategije generiranja ovisno o omjeru
sigurnosnih razina originalnog primjera $sl_p$ i susjeda $sl_n$: ako su
obje niske, generiranje se izbjegava; ako je jedna nula, koristi se kopija
sigurnijeg primjera; ako je $sl_p < sl_n$, novi primjer smje\v{s}ta se
bli\v{z}e sigurnijem. Dat \'{c}e se pseudo-k\^{o}d. Izbjegavanje
generiranja u nesigurnim podru\v{c}jima je prednost, dok je nedostatak
konzervativnost koja mo\v{z}e ograni\v{c}iti pro\v{s}irenje manjinske
klase.

\subsection{K-Means SMOTE}

K-Means SMOTE \cite{douzas2018} rje\v{s}ava problem unutar-klasnog
disbalansa. Postupak: prvo se manjinski primjeri grupiraju K-Means
algoritmom, zatim se identificiraju rijetki klasteri i dodjeljuje im se
proporcionalno vi\v{s}e sinteti\v{c}kih primjera (te\v{z}ina $\propto
1/\text{veli\v{c}ina\_klastera}$), nakon \v{c}ega se SMOTE primjenjuje
unutar svakog klastera koriste\'{c}i samo susjede iz istog klastera.
Dat \'{c}e se pseudo-k\^{o}d. Istovremeno rje\v{s}avanje me\dj{}u-klasnog
i unutar-klasnog disbalansa predstavlja glavnu prednost, dok su dodatni
parametar (broj klastera) i ve\'{c}a ra\v{c}unska zahtjevnost glavni
nedostaci. Srodni pristupi temeljeni na filtriranju obra\dj{}eni su
u radu \cite{saez2015}.

\subsection{SVM-SMOTE}

SVM-SMOTE \cite{nguyen2011} koristi SVM klasifikator za identifikaciju
potpornih vektora manjinske klase koji le\v{z}e na granici odluke.
Postupak: trenira se SVM na originalnim podacima, identificiraju se
potporni vektori manjinske klase, a zatim se generiraju sinteti\v{c}ki
primjeri du\v{z} linija koje spajaju te potporne vektore s njihovim
najbli\v{z}im susjedima. Dat \'{c}e se pseudo-k\^{o}d. Precizno
pozicioniranje novih primjera na granici odluke je klju\v{c}na prednost,
a ve\'{c}a ra\v{c}unska slo\v{z}enost zbog treniranja SVM-a glavni
nedostatak.

\subsection{SMOTE-ENN i SMOTE-Tomek}

SMOTE-ENN i SMOTE-Tomek \cite{batista2004} kombiniraju preuzorkovanje s
naknadnim \v{c}i\v{s}\'{c}enjem skupa podataka. SMOTE-ENN prvo generira
sinteti\v{c}ke primjere, a zatim uklanja sve primjere koje ve\'{c}ina
$k$-NN susjeda pogre\v{s}no klasificira. SMOTE-Tomek tako\dj{}er prvo
primjenjuje SMOTE, a zatim uklanja Tomek Link parove --- parove primjera
razli\v{c}itih klasa koji su me\dj{}usobno najbli\v{z}i. Dat \'{c}e se
pseudo-k\^{o}d za oba pristupa. Prednost je smanjenje \v{s}uma i
preklapanja klasa, dok je nedostatak mogu\'{c}nost gubitka korisnih
primjera.

\subsection{Geometrijska pro\v{s}irenja SMOTE-a}

Izvorni SMOTE generira sinteti\v{c}ke primjere isklju\v{c}ivo na linijskom
segmentu izme\dj{}u dva postoje\'{c}a primjera, \v{s}to ograni\v{c}ava
raznolikost generiranih uzoraka. U nastavku se opisuju tri geometrijska
pro\v{s}irenja koja adresiraju ovo ograni\v{c}enje.

G-SMOTE (Geometric SMOTE) \cite{douzas2019} generira sinteti\v{c}ke
primjere unutar vi\v{s}edimenzionalnog geometrijskog sektora definiranog
manjinskim primjerom kao sredi\v{s}tem i smjerom prema odabranom susjedu.
Uvodi dva dodatna parametra: faktor deformacije i faktor skra\'{c}ivanja
koji kontroliraju koliko daleko od linije se primjer mo\v{z}e generirati.
G-SMOTE se mo\v{z}e koristiti kao izravna zamjena za osnovni SMOTE.

Random SMOTE \cite{dong2011} za svaki manjinski primjer nasumi\v{c}no bira
smjer i udaljenost za generiranje, \v{c}ime pove\'{c}ava raznolikost ali
mo\v{z}e generirati u podru\v{c}jima ve\'{c}inske klase.

Polynom-Fit SMOTE \cite{gazzah2008} koristi polinomnu interpolaciju kroz
vi\v{s}e susjeda umjesto linearnog segmenta, \v{s}to bolje prati stvarnu
geometriju podataka kod nelinearnih distribucija. Jo\v{s} jedno
zna\v{c}ajno geometrijsko/te\v{z}insko pro\v{s}irenje koje nije
implementirano u ovom radu je MWMOTE \cite{barua2014}.

\subsection{Usporedni pregled izvedenica}

Dat \'{c}e se tablica sa sumarnim pregledom svih obra\dj{}enih izvedenica,
uklju\v{c}uju\'{c}i referencu, godinu objave, klju\v{c}ni problem koji
rje\v{s}ava, osnovni koncept i klju\v{c}ne parametre. Dodatne varijante
poput MWMOTE \cite{barua2014} i SMOTE-IPF \cite{saez2015} postoje u
literaturi, no izvan su opsega ove eksperimentalne usporedbe.

\subsection{Alternativni pristupi rje\v{s}avanju problema}

\subsubsection{Metode poduzorkovanja}

Opisat \'{c}e se metode koje smanjuju dominaciju ve\'{c}inske klase
uklanjanjem njenih primjera: NearMiss, Tomek Links, Edited Nearest
Neighbors (ENN) i One-Sided Selection (OSS). Za svaku \'{c}e se objasniti
princip rada i situacije u kojima se koristi. Koristit \'{c}e se
literatura \cite{batista2004,liu2009}.

\subsubsection{Cost-sensitive u\v{c}enje}

Opisat \'{c}e se princip cost-sensitive u\v{c}enja \cite{elkan2001} kod
kojeg se razli\v{c}itim klasama dodjeljuju razli\v{c}ite cijene pogre\v{s}ne
klasifikacije, \v{c}ime se klasifikator prisiljava da posveti vi\v{s}e
pa\v{z}nje manjinskoj klasi. Objasnit \'{c}e se implementacija putem
matrice tro\v{s}kova i class\_weight parametra. Navest \'{c}e se prednosti
(ne mijenja distribuciju podataka) i nedostaci (zahtijeva pa\v{z}ljivo
definiranje tro\v{s}kova).

\subsubsection{Ansambl metode}

Opisat \'{c}e se ansambl pristupi koji kombiniraju resampliranje s
bagging i boosting metodama: SMOTEBoost \cite{chawla2003} primjenjuje
SMOTE u svakoj iteraciji AdaBoost-a, RUSBoost \cite{seiffert2010} koristi
nasumi\v{c}no poduzorkovanje, EasyEnsemble i BalanceCascade \cite{liu2009}
koriste vi\v{s}estruko poduzorkovanje ve\'{c}inske klase, a
BalancedRandomForest balansira bootstrap uzorke. Pregled \'{c}e se
temeljiti na radu \cite{galar2012}.

\subsubsection{Pristupi zasnovani na neuronskim mre\v{z}ama}

Opisat \'{c}e se Focal Loss funkcija \cite{lin2017} koja modificira
cross-entropy smanjuju\'{c}i doprinos lakih primjera, te generativne
suprotstavljene mre\v{z}e \cite{goodfellow2014} za generiranje
visokokvalitetnih sinteti\v{c}kih primjera manjinske klase. Navest
\'{c}e se prednosti (kvaliteta generiranja) i nedostaci (veliki zahtjevi
za koli\v{c}inom podataka i ra\v{c}unskim resursima). Ovaj dio \'{c}e biti
pregledan jer fokus rada ostaje na SMOTE algoritmima.
